% Two-page student handout, DSCT Week 13 Tuesday (Nov 10, 2026). Values checked by presentations/beamer/week13_tue_boolean_binary/check_numbers.py
\documentclass[10pt,letterpaper]{article}
\usepackage[margin=0.6in]{geometry}
\usepackage{amssymb,amsmath}
\usepackage{array}\usepackage{xcolor}
\pagestyle{empty}
\setlength{\parindent}{0pt}\setlength{\parskip}{3pt}
\renewcommand{\arraystretch}{1.05}
\begin{document}
{\Large\textbf{Counting, Carrying, and Circuits}}\hfill COMSC 2043 \ \textbullet\ \ Nov 10, 2026\\
\rule{\textwidth}{0.6pt}

\textbf{1. Carrying is the same in every base.} A base-$b$ system has $b$ symbols; the biggest is $b-1$. When a place runs out of symbols, write 0 there and add 1 to the next place.
\begin{center}
\begin{tabular}{l|ccc}
& \textbf{biggest digit $+1$} & \textbf{two biggest $+1$} & \textbf{place values}\\\hline
base 10 & $9+1=10$ & $99+1=100$ & $1,10,100$\\
base 8 & $7+1=10$ & $77+1=100$ & $1,8,64$\\
base 16 & $F+1=10$ & $FF+1=100$ & $1,16,256$\\
base 2 & $1+1=10$ & $11+1=100$ & $1,2,4,8$
\end{tabular}
\end{center}
``10'' always means \emph{one group of $b$, none left over}. The marks could be anything: circle, star, face, fish.

\textbf{2. Add one column at a time, right to left, then move sideways.} Line the digits up. Write any carry above the next column.
\newcommand{\gc}[1]{\textcolor{black!60}{\footnotesize #1}}
\begin{center}
\renewcommand{\arraystretch}{1.15}
{\Large\begin{tabular}{@{}l@{\ \ }c@{\ \ }c@{\ \ }c@{}}
\gc{place} & \gc{100s} & \gc{10s} & \gc{1s}\\
\gc{carry} & \gc{1} & & \\
 & & \texttt{5} & \texttt{2}\\
$+$ & & \texttt{5} & \texttt{3}\\\hline
 & \texttt{1} & \texttt{0} & \texttt{5}
\end{tabular}}
\qquad
\begin{minipage}[c]{0.36\textwidth}\small
\textbf{Ones:} $2+3=5$.\\
\textbf{Tens:} $5+5=10$. Write 0, \textbf{carry 1}.\\
\textbf{100s:} nothing to add, so the carry drops down: 1.\\
\textbf{Answer: 105.}
\end{minipage}
\end{center}

\textbf{Binary uses the identical steps}, with a smaller table: $0+0=0,\ 0+1=1,\ 1+0=1,\ 1+1=\texttt{10}$ (write 0, carry 1). Each example also shows the numbers converted to decimal, so you can check that the math works the same way \emph{both directions}: binary $\to$ decimal, and decimal $\to$ binary.

\textbf{Short} \ ($3+1=4$)\\[2pt]
{\large\begin{tabular}{@{}l>{\centering\arraybackslash}p{0.78cm}>{\centering\arraybackslash}p{0.78cm}>{\centering\arraybackslash}p{0.78cm}@{}}
\gc{place} & \gc{4s} & \gc{2s} & \gc{1s}\\
\gc{carry} & \gc{1} & \gc{1} & \\
 &  & \texttt{1} & \texttt{1}\\
$+$ &  &  & \texttt{1}\\\hline
 & \texttt{1} & \texttt{0} & \texttt{0}
\end{tabular}}\\[1pt]
{\small \texttt{11} $=2+1=3$; \ \texttt{1} $=1=1$; \ sum $3+1=4=4=$ \texttt{100}}

\textbf{Medium} \ ($11+6=17$)\\[2pt]
\begin{minipage}[c]{0.5\textwidth}{\large\begin{tabular}{@{}l>{\centering\arraybackslash}p{0.78cm}>{\centering\arraybackslash}p{0.78cm}>{\centering\arraybackslash}p{0.78cm}>{\centering\arraybackslash}p{0.78cm}>{\centering\arraybackslash}p{0.78cm}@{}}
\gc{place} & \gc{16s} & \gc{8s} & \gc{4s} & \gc{2s} & \gc{1s}\\
\gc{carry} & \gc{1} & \gc{1} & \gc{1} & \gc{0} & \\
 &  & \texttt{1} & \texttt{0} & \texttt{1} & \texttt{1}\\
$+$ &  & \texttt{0} & \texttt{1} & \texttt{1} & \texttt{0}\\\hline
 & \texttt{1} & \texttt{0} & \texttt{0} & \texttt{0} & \texttt{1}
\end{tabular}}\end{minipage}
\begin{minipage}[c]{0.45\textwidth}\small
\textbf{1s:} $1+0=1$. \ \textbf{2s:} $1+1=10$, write 0, carry 1.\\
\textbf{4s:} $0+1+1=10$, write 0, carry 1.\\
\textbf{8s:} $1+0+1=10$, write 0, carry 1.\\
\textbf{16s:} the carry drops down: 1.
\end{minipage}\\[1pt]
{\small \texttt{1011} $=8+2+1=11$; \ \texttt{0110} $=4+2=6$; \ sum $11+6=17=16+1=$ \texttt{10001}}

\textbf{Long} \ ($183+109=292$)\\[2pt]
{\large\begin{tabular}{@{}l>{\centering\arraybackslash}p{0.78cm}>{\centering\arraybackslash}p{0.78cm}>{\centering\arraybackslash}p{0.78cm}>{\centering\arraybackslash}p{0.78cm}>{\centering\arraybackslash}p{0.78cm}>{\centering\arraybackslash}p{0.78cm}>{\centering\arraybackslash}p{0.78cm}>{\centering\arraybackslash}p{0.78cm}>{\centering\arraybackslash}p{0.78cm}@{}}
\gc{place} & \gc{256s} & \gc{128s} & \gc{64s} & \gc{32s} & \gc{16s} & \gc{8s} & \gc{4s} & \gc{2s} & \gc{1s}\\
\gc{carry} & \gc{1} & \gc{1} & \gc{1} & \gc{1} & \gc{1} & \gc{1} & \gc{1} & \gc{1} & \\
 &  & \texttt{1} & \texttt{0} & \texttt{1} & \texttt{1} & \texttt{0} & \texttt{1} & \texttt{1} & \texttt{1}\\
$+$ &  & \texttt{0} & \texttt{1} & \texttt{1} & \texttt{0} & \texttt{1} & \texttt{1} & \texttt{0} & \texttt{1}\\\hline
 & \texttt{1} & \texttt{0} & \texttt{0} & \texttt{1} & \texttt{0} & \texttt{0} & \texttt{1} & \texttt{0} & \texttt{0}
\end{tabular}}\\[1pt]
{\small \texttt{10110111} $=128+32+16+4+2+1=183$; \ \texttt{01101101} $=64+32+8+4+1=109$; \ sum $183+109=292=256+32+4=$ \texttt{100100100}}
\vspace{2pt}
\newpage
\textbf{3. Multiply} the same way in every base: one partial product per digit of the bottom number, each shifted left one place, then add the columns. In binary a partial product is either a copy of the top number or all zeros.

\textbf{Decimal first} \ ($23\times14=322$)\\[2pt]
\begin{minipage}[c]{0.3\textwidth}
{\large\begin{tabular}{@{}r>{\centering\arraybackslash}p{0.78cm}>{\centering\arraybackslash}p{0.78cm}>{\centering\arraybackslash}p{0.78cm}@{}}
 & & \texttt{2} & \texttt{3}\\
$\times$ & & \texttt{1} & \texttt{4}\\\hline
 & & \texttt{9} & \texttt{2}\\
 & \texttt{2} & \texttt{3} & \texttt{0}\\\hline
 & \texttt{3} & \texttt{2} & \texttt{2}
\end{tabular}}
\end{minipage}
\begin{minipage}[c]{0.6\textwidth}\small
$23\times4=92$ \quad (the ones digit of the bottom number)\\
$23\times10=230$ \quad (the tens digit: shifted left one place)\\
$92+230=322$. Add the columns, right to left, with carries.
\end{minipage}

\textbf{Binary example 1} \ ($5\times3=15$)\\[2pt]
{\large\begin{tabular}{@{}r>{\centering\arraybackslash}p{0.78cm}>{\centering\arraybackslash}p{0.78cm}>{\centering\arraybackslash}p{0.78cm}>{\centering\arraybackslash}p{0.78cm}@{}}
 &  & \texttt{1} & \texttt{0} & \texttt{1}\\
$\times$ &  &  & \texttt{1} & \texttt{1}\\\hline
 &  & \texttt{1} & \texttt{0} & \texttt{1}\hfill\gc{}\\
 & \texttt{1} & \texttt{0} & \texttt{1} & \hfill\gc{}\\\hline
 & \texttt{1} & \texttt{1} & \texttt{1} & \texttt{1}
\end{tabular}}\\[1pt]
{\small \texttt{101} $=5$; \ \texttt{11} $=3$; \ product $5\times3=15=8+4+2+1=$ \texttt{1111}}

\textbf{Binary example 2} \ ($6\times5=30$)\\[2pt]
{\large\begin{tabular}{@{}r>{\centering\arraybackslash}p{0.78cm}>{\centering\arraybackslash}p{0.78cm}>{\centering\arraybackslash}p{0.78cm}>{\centering\arraybackslash}p{0.78cm}>{\centering\arraybackslash}p{0.78cm}@{}}
 &  &  & \texttt{1} & \texttt{1} & \texttt{0}\\
$\times$ &  &  & \texttt{1} & \texttt{0} & \texttt{1}\\\hline
 &  &  & \texttt{1} & \texttt{1} & \texttt{0}\hfill\gc{}\\
 &  & \texttt{0} & \texttt{0} & \texttt{0} & \hfill\gc{}\\
 & \texttt{1} & \texttt{1} & \texttt{0} &  & \hfill\gc{}\\\hline
 & \texttt{1} & \texttt{1} & \texttt{1} & \texttt{1} & \texttt{0}
\end{tabular}}\\[1pt]
{\small \texttt{110} $=6$; \ \texttt{101} $=5$; \ product $6\times5=30=16+8+4+2=$ \texttt{11110}}

\vspace{2pt}\rule{\textwidth}{0.3pt}

\textbf{4. Boolean algebra.} Values 0 and 1. \textbf{AND} $A\cdot B$ is 1 only if both are 1. \textbf{OR} $A+B$ is 1 if at least one is 1. \textbf{NOT} $\overline{A}$ flips. \textbf{XOR} $A\oplus B$ is 1 when the inputs differ.

\textbf{Careful:} in Boolean algebra $1+1=1$ (OR). In binary arithmetic $1+1=10$ (add). Always say which system you are in.

\textbf{5. Half adder} (adds two bits). Sum $S=A\oplus B$. Carry $C=A\cdot B$. Read $CS$ as a two-bit number: it equals $A+B$.
\begin{center}
\begin{tabular}{cc|cc|c}
$A$&$B$&$C$&$S$&$A+B$ in binary\\\hline
0&0&0&0&\texttt{00}\\
0&1&0&1&\texttt{01}\\
1&0&0&1&\texttt{01}\\
1&1&1&0&\texttt{10}
\end{tabular}
\end{center}

\textbf{6. Full adder} (adds $A$, $B$, and a carry-in). $S=A\oplus B\oplus C_{in}$ \quad $C_{out}=A\cdot B+C_{in}\cdot(A\oplus B)$.
\begin{center}
\begin{tabular}{ccc|cc|c}
$A$&$B$&$C_{in}$&$C_{out}$&$S$&$A+B+C_{in}$\\\hline
0&0&0&0&0&\texttt{00}\\
0&0&1&0&1&\texttt{01}\\
0&1&0&0&1&\texttt{01}\\
0&1&1&1&0&\texttt{10}\\
1&0&0&0&1&\texttt{01}\\
1&0&1&1&0&\texttt{10}\\
1&1&0&1&0&\texttt{10}\\
1&1&1&1&1&\texttt{11}
\end{tabular}
\end{center}
All eight input rows are listed, so the table is a \textbf{proof}, not just an example.

\textbf{7. Chaining adders.} One full adder per bit, each carry feeding the next: a \textbf{ripple-carry adder}. The carry travels through every stage, so about $2n$ gate delays for $n$ bits (8 bits: 16; 16 bits: 32). \textbf{Carry lookahead} computes the carries directly from \emph{generate} $G_i=A_iB_i$ and \emph{propagate} $P_i=A_i\oplus B_i$: $C_1=G_0+P_0C_0$, $C_2=G_1+P_1G_0+P_1P_0C_0$. A 16-bit lookahead adder finishes in about 8 gate delays instead of 32, for more hardware.
\end{document}
